\documentclass[12pt,a4paper]{ctexart}

% 页面设置
\usepackage[margin=2.5cm,headheight=14pt]{geometry}
\usepackage{setspace}
\onehalfspacing

% 颜色与样式
\usepackage{xcolor}
\definecolor{evogreen}{RGB}{34,139,34}
\definecolor{evoyellow}{RGB}{218,165,32}
\definecolor{evored}{RGB}{178,34,34}
\definecolor{evogray}{RGB}{100,100,100}
\definecolor{codebg}{RGB}{248,248,248}
\definecolor{linkblue}{RGB}{30,90,150}

\usepackage{hyperref}
\hypersetup{
  colorlinks=true,
  linkcolor=linkblue,
  urlcolor=linkblue,
  citecolor=linkblue,
  pdftitle={Evo-Kernel 用户手册},
  pdfauthor={zodyne}
}

% 图表
\usepackage{tikz}
\usetikzlibrary{shapes.geometric, arrows.meta, positioning, fit, backgrounds, calc, shadows}
\usepackage{booktabs}
\usepackage{graphicx}
\usepackage{float}

% 代码块
\usepackage{listings}
\lstset{
  backgroundcolor=\color{codebg},
  basicstyle=\ttfamily\small,
  breaklines=true,
  frame=single,
  rulecolor=\color{evogray!30},
  keywordstyle=\color{linkblue}\bfseries,
  commentstyle=\color{evogray},
  stringstyle=\color{evogreen},
  showstringspaces=false,
  xleftmargin=1em,
  xrightmargin=1em,
  aboveskip=1em,
  belowskip=1em
}

% 提示框
\usepackage{tcolorbox}
\tcbuselibrary{skins,breakable}
\newtcolorbox{tipbox}[1][]{
  colback=evogreen!5,
  colframe=evogreen!70,
  fonttitle=\bfseries,
  title=提示,
  #1
}
\newtcolorbox{warnbox}[1][]{
  colback=evoyellow!5,
  colframe=evoyellow!80!black,
  fonttitle=\bfseries,
  title=注意,
  #1
}
\newtcolorbox{redbox}[1][]{
  colback=evored!5,
  colframe=evored!80!black,
  fonttitle=\bfseries,
  title=红线,
  #1
}

% 页眉页脚
\usepackage{fancyhdr}
\pagestyle{fancy}
\fancyhf{}
\fancyhead[L]{\small Evo-Kernel 用户手册}
\fancyhead[R]{\small \nouppercase{\leftmark}}
\fancyfoot[C]{\thepage}
\renewcommand{\headrulewidth}{0.4pt}

% 标题
\title{
  \vspace{-2cm}
  \textbf{Evo-Kernel}\\[0.3em]
  \Large 个人经验治理与固化层\\[0.2em]
  \large 使用手册
}
\author{zodyne}
\date{\today}

\begin{document}

\maketitle
\thispagestyle{empty}

\begin{abstract}
\noindent
Evo-Kernel 是位于通用记忆框架之上的「经验治理与固化层」。它以纯文件 + git 作为单一事实源，通过捕获、蒸馏、人审、固化四个环节，把零散的工作经验转化为可检索、可审计、可硬化的约束与技能。

本手册面向日常使用者，覆盖：快速开始、核心工作流、周期维护、备份恢复、以及使用中的红线与注意事项。
\end{abstract}

\tableofcontents
\newpage

%======================================================================
\section{Evo-Kernel 是什么}
%======================================================================

\subsection{一句话定位}

通用记忆框架（Mem0 / Zep / Letta 等）解决的是经验的「存储与检索」；Evo-Kernel 解决的是其上没人解决的「治理、固化、可审计」问题。二者是叠加而非替代。

\subsection{核心能力}

\begin{itemize}
  \item \textbf{捕获}：零判断入口，任何想法先丢进 inbox。
  \item \textbf{检索注入}：按任务自动从 playbook / facts / episodes / principles 召回相关经验。
  \item \textbf{离线蒸馏}：session 结束后对 transcript 切片，由 Reflector 生成提案，人审后入库。
  \item \textbf{技能固化}：经验 $\rightarrow$ SKILL.md，可直接被 Claude Code / Pi 挂载。
  \item \textbf{硬约束硬化}：经验 $\rightarrow$ guard / hook，对危险工具调用进行 warn / block。
  \item \textbf{git 可审计}：每一次 curate / solidify / demote 都是一次可 diff、可 revert 的 commit。
\end{itemize}

\subsection{架构分层}

图 \ref{fig:architecture} 展示了 Evo-Kernel 的层次结构。上层是 harness（Claude Code / Pi），下层是纯文件 + git 的内核，中间通过 hooks 与 CLI 交互。

\begin{figure}[H]
\centering
\begin{tikzpicture}[
  box/.style={rectangle, rounded corners, draw, minimum width=3.2cm, minimum height=0.9cm, align=center, font=\small},
  layer/.style={rectangle, rounded corners, draw=evogray, dashed, inner sep=0.4cm},
  arrow/.style={-{Stealth[length=2mm]}, thick},
  label/.style={font=\footnotesize\bfseries, color=evogray}
]

% Harness 层
\node[box, fill=evogreen!15] (cc) {Claude Code};
\node[box, fill=evogreen!15, right=1.2cm of cc] (pi) {Pi};

% 适配层
\node[box, fill=evoyellow!15, below=1.2cm of cc] (hook1) {UserPromptSubmit\\SessionEnd\\PreToolUse};
\node[box, fill=evoyellow!15, below=1.2cm of pi] (hook2) {before\_agent\_start\\session\_shutdown\\tool\_call};

% 协议层
\node[box, fill=linkblue!10, below=1.5cm of hook1] (proto) {evo CLI\\26 个命令};

% 内核层
\node[box, fill=evored!10, below=1.2cm of proto] (kernel) {内核层\\目录即状态机};

% 外挂层
\node[box, fill=evogray!15, right=2.5cm of kernel] (ext) {可选外挂\\Mem0 / Cognee / LightRAG};

% 箭头
\draw[arrow] (cc) -- (hook1);
\draw[arrow] (pi) -- (hook2);
\draw[arrow] (hook1) -- (proto);
\draw[arrow] (hook2) -- (proto);
\draw[arrow] (proto) -- (kernel);
\draw[arrow, dashed, color=evogray] (kernel) -- (ext) node[midway, above, font=\scriptsize] {MCP / opt-in};

\end{tikzpicture}
\caption{Evo-Kernel 架构分层}
\label{fig:architecture}
\end{figure}

%======================================================================
\section{目录与状态机}
%======================================================================

\subsection{目录角色}

Evo-Kernel 的每个目录都有明确的角色，这是「目录即状态机」的体现：文件在不同目录间移动，即代表经验状态的变化。

\begin{table}[H]
\centering
\caption{目录角色与注入权限}
\begin{tabular}{@{}llc@{}}
\toprule
\textbf{目录} & \textbf{角色} & \textbf{是否注入} \\
\midrule
\texttt{inbox/} & capture 暂存 + 会话登记 & \textcolor{evored}{否} \\
\texttt{lessons/} & candidate 经验暂存 & \textcolor{evored}{否} \\
\texttt{facts/} & 事实/偏好/环境 & \textcolor{evogreen}{是} \\
\texttt{episodes/} & 情景记忆：一次任务一份 & \textcolor{evogreen}{是} \\
\texttt{playbook/} & 策略库，带 helpful/harmful 计数 & \textcolor{evogreen}{是} \\
\texttt{principles/} & 跨域普适原则 & \textcolor{evogreen}{是} \\
\texttt{skills/} & SKILL.md 包（agentskills.io） & 经 skill 注入 \\
\texttt{ops/archive/} & 退役/固化存档 & \textcolor{evored}{否} \\
\texttt{ops/constraints/} & 硬约束规则（JSON） & guard 执行 \\
\texttt{ops/log/} & 运行日志（全部 \texttt{gitignore}） & — \\
\texttt{ops/proposals/} & reflect/distill 提案 & 待人审 \\
\texttt{index/manifest.yaml} & 派生索引（可重建） & — \\
\bottomrule
\end{tabular}
\end{table}

\begin{redbox}
\textbf{不变量 I2：} \texttt{inbox/} 与 \texttt{lessons/} 中的条目永远不会被自动注入 context。只有经过人审 curate 进入 playbook / facts / episodes / principles 的条目，才有资格被 recall。否则 capture 这个零审查入口就成了绕过人审的后门。
\end{redbox}

\subsection{经验生命周期}

图 \ref{fig:lifecycle} 展示了一条经验从捕获到固化或退役的完整路径。

\begin{figure}[H]
\centering
% 布局：主链路（捕获→提案→入库）横向在 y=0，两条固化分支上下分岔，退役在下方。
% superseded 回边走右侧页边（x=14.3）绕开 hook 节点，不穿过节点场。
% 宽度核算（textwidth = 16cm）：节点半宽 ≤1.41，横向间距 4.3/4.4 → 节点间空档 1.7cm，
% 恰好容下最长边标签 "slice / reflect"（\scriptsize ≈1.5cm）；
% 总占位 = 左缘 -1.3 至右缘 14.3 = 15.6cm，留 0.4cm 余量。
\begin{tikzpicture}[
  state/.style={rectangle, rounded corners, draw, minimum width=2.6cm, minimum height=1.0cm, align=center, font=\footnotesize},
  arrow/.style={-{Stealth[length=2mm]}, thick},
  annot/.style={font=\scriptsize, color=evogray, fill=white, inner sep=1pt}
]

\node[state, fill=evogreen!10]  (capture) at (0,0)      {inbox\\捕获};
\node[state, fill=evoyellow!10] (slice)   at (4.3,0)    {ops/proposals\\提案};
\node[state, fill=linkblue!10]  (curate)  at (8.7,0)    {playbook/facts/…\\入库};
\node[state, fill=evored!10]    (skill)   at (12.7,1.9) {skills/\\SKILL 固化};
\node[state, fill=evored!10]    (hook)    at (12.7,-1.9){ops/constraints/\\hook 硬化};
\node[state, fill=evogray!20]   (archive) at (8.7,-3.6) {ops/archive/\\退役};

\draw[arrow] (capture) -- (slice);
\draw[arrow] (slice) -- (curate);
\draw[arrow] (curate) -- (skill);
\draw[arrow] (curate) -- (hook);
\draw[arrow] (curate) -- (archive);
\draw[arrow, dashed] (skill.east) -- (14.3,1.9) -- (14.3,-3.6) -- (archive.east);

% 边标签放在节点上沿之上（y=0.78 > 节点半高 0.5），横向再宽也不会压到节点；
% 两条斜向标签退到 x=10.4，避开 skill/hook 节点左缘（x≈11.4）。
\node[annot] at (2.15,0.78)  {slice / reflect};
\node[annot] at (6.5,0.78)   {人审 curate};
\node[annot] at (10.15,1.2)  {solidify skill};
\node[annot] at (10.15,-1.2) {solidify hook};
\node[annot, anchor=west] at (8.85,-1.9) {demote};
\node[annot] at (12.2,-3.6)  {superseded};

\end{tikzpicture}
\caption{经验生命周期状态机}
\label{fig:lifecycle}
\end{figure}

%======================================================================
\section{快速开始}
%======================================================================

\subsection{日常检查命令}

每天开始工作前，先检查系统健康：

\begin{lstlisting}[language=bash]
cd ~/Dev/evo-kernel
./bin/evo doctor
\end{lstlisting}

期望输出是 \texttt{OK (N PASS / 少量 WARN / 0 FAIL)}，只要 FAIL=0 即为健康。常见的 WARN 有三类，均非故障：工作区有未提交变更、manifest 滞后（跑一次 \texttt{evo index rebuild} 即消）、transcript 哨兵比（已知飞轮状态，由 reflect 周期跟踪）。

\subsection{捕获一个想法}

任何时候想到一个值得记录的经验、坑或偏好，直接 capture：

\begin{lstlisting}[language=bash]
./bin/evo capture "遇到 Claude Code 权限拦截时，先用 AskUserQuestion 或明确授权，不要偷偷绕过"
\end{lstlisting}

这会在 \texttt{inbox/} 下生成一个带时间戳的 markdown 文件。capture 是零判断入口，不需要结构化，不需要 frontmatter。

\begin{tipbox}
capture 的目标不是写完美条目，而是不错过任何一闪而过的经验。结构化和人审留给后面的 curate。
\end{tipbox}

\subsection{主动召回经验}

如果你在写代码或调试，想主动看看有没有相关经验：

\begin{lstlisting}[language=bash]
./bin/evo recall --task "如何安全地删除远程 git 分支"
\end{lstlisting}

这会输出一段可注入的文本（如果命中）。更常见的场景是 Claude Code / Pi 自动在每次输入时调用 \texttt{hook-recall}，你不需要手动跑。

%======================================================================
\section{核心工作流}
%======================================================================

\subsection{工作流一：从 capture 到 curate}

这是最常见的单条经验入库路径。

\textbf{Step 1：把 inbox 里的 raw 文本整理成提案}

\begin{lstlisting}[language=bash]
./bin/evo inbox
# 看到一条 capture，打开编辑
# 把它改写成符合 SCHEMA.md 的提案文件，放到 ops/proposals/
\end{lstlisting}

一个最小提案示例：

\begin{lstlisting}[language=Markdown]
---
id: avoid-silent-bypass
name: 权限拦截不绕过
type: principle
status: candidate
scope: [workflow]
domains: [claude-code]
triggers:
  - 权限拦截
  - Claude Code 拒绝
  - bypass
evidence:
  helpful: 0
  harmful: 0
verified_by: human
---

当 Claude Code 的权限系统拦截一个操作（如删除远程分支、修改系统配置）时，
正确动作是：
1. 用 AskUserQuestion 向用户明确请求授权；
2. 或把命令贴给用户，让用户用 `!` 前缀自己执行。

错误动作是：
- 用其他工具绕过权限拦截；
- 把破坏性操作伪装成无害命令。
\end{lstlisting}

\textbf{Step 2：人审并入库}

\begin{lstlisting}[language=bash]
./bin/evo curate --file ops/proposals/avoid-silent-bypass.md --to principles
\end{lstlisting}

curate 会做三件事：
\begin{enumerate}
  \item schema 校验（id / type / status / triggers 必填）。
  \item 敏感信息检查（凭据、密钥、内网地址、PII 不得入库）。
  \item 指令样内容审查（正文必须是陈述式经验，不能是面向未来 agent 的祈使指令）。
\end{enumerate}

通过后，文件被移动到 \texttt{principles/}，自动 rebuild manifest、commit、并 push 到 remote。

\begin{warnbox}
curate 是唯一的入库口（不变量 I3）。不要手动把文件拖进 \texttt{playbook/} 或 \texttt{facts/}，否则绕过了敏感信息和指令样审查。
\end{warnbox}

\subsection{工作流二：会话蒸馏}

一次完整工作会话结束后，把 transcript 中的经验提炼出来。

\begin{tipbox}
这条流程现在已经自动化：SessionEnd 触发器会调用 \texttt{ops/bin/evo-distill.sh}（带锁，防并发），launchd 每日 03:30 再跑一次兜底（\texttt{--max 3}），清掉会话结束时机器休眠、harness 超时或锁冲突漏掉的积压。队列由 \texttt{evo queue} 给出、完成后 \texttt{evo mark-distilled} 回写。下面的手动步骤保留，用于自动链路失效时的排查与补跑。
\end{tipbox}

\textbf{Step 1：确认 session 已登记}

Claude Code / Pi 的 session-end hook 会自动写 \texttt{inbox/session-refs.jsonl}。你也可以手动登记：

\begin{lstlisting}[language=bash]
./bin/evo session-end --session /path/to/transcript.jsonl --id <session-id>
\end{lstlisting}

\textbf{Step 2：切片}

\begin{lstlisting}[language=bash]
./bin/evo slice --session /path/to/transcript.jsonl --ids <session-id>
\end{lstlisting}

这会输出一个 Markdown 报告，包含：命令 $\leftrightarrow$ 结果对照、写/改文件列表、首条 user、末条 assistant，以及该 session 被注入了哪些经验条目。

\textbf{Step 3：生成 reflect 提案}

\begin{lstlisting}[language=bash]
./bin/evo reflect --save
\end{lstlisting}

reflect 会：
\begin{itemize}
  \item 统计库状态（按 zone、status 分桶）；
  \item 分析 recall.jsonl（调用数、命中数、空命中任务）；
  \item 跑 audit，列出治理问题；
  \item 输出判据对照表（M1/M2/P4 判据、transcript 时效、降级盘点等）；
  \item \texttt{--save} 时把报告写入 \texttt{ops/proposals/reflect-<日期>.md}。
\end{itemize}

\textbf{Step 4：人审并 curate}

打开 \texttt{ops/proposals/reflect-<日期>.md}，把认可的经验改写成标准提案文件，然后：

\begin{lstlisting}[language=bash]
./bin/evo curate --file ops/proposals/<proposal>.md --to <playbook|facts|episodes|principles>
\end{lstlisting}

\subsection{工作流三：技能固化}

当一条经验反复有用、值得变成可挂载的技能时：

\begin{lstlisting}[language=bash]
./bin/evo solidify --id avoid-silent-bypass --to skill
\end{lstlisting}

这会：
\begin{enumerate}
  \item 在 \texttt{skills/avoid-silent-bypass/SKILL.md} 创建技能包；
  \item 把原条目标记为 \texttt{solidified}，写入 \texttt{superseded\_by: skill:avoid-silent-bypass}，移到 \texttt{ops/archive/}；
  \item 自动 rebuild manifest、调用 \texttt{evo link} 同步软链、commit、push。
\end{enumerate}

之后，该 skill 会出现在 \texttt{\textasciitilde/.claude/skills/} 软链中，Claude Code 可以加载。

\subsection{工作流四：批量导入外部资料与技术资料}

如果你已经有大量技术笔记、文档摘录、会议记录或项目沉淀（例如从 Obsidian、Notion、Markdown 仓库导出），不要直接拖进 \texttt{facts/} 或 \texttt{playbook/}——那样会绕过 curate 的人审、敏感信息检查和指令样审查（违反不变量 I3）。

\subsubsection{导入的两阶段模型}

借鉴 Zettelkasten 的 ``fleeting note → literature note → permanent note'' 流程，以及 RAG ingestion pipeline 的 ``extract → clean → chunk → enrich → index'' 模式，Evo-Kernel 推荐采用 \textbf{机器预处理 + 人工审查} 的两阶段模型：

\begin{figure}[H]
\centering
% 布局：4 阶段线性流水线，绝对坐标 x = 0 / 4.2 / 8.4 / 12.6。
% 宽度核算（textwidth = 16cm）：最宽节点内容 ~2.2cm + inner sep 0.67cm = 2.87cm，
% 半宽 1.44 → 横向占位 = 12.6 + 2×1.44 = 15.5cm，留 0.5cm 余量。
\begin{tikzpicture}[
  stage/.style={rectangle, rounded corners, draw, minimum width=2.6cm, minimum height=1.1cm, align=center, font=\footnotesize},
  arrow/.style={-{Stealth[length=2mm]}, thick},
  annot/.style={font=\scriptsize, color=evogray, fill=white, inner sep=1pt}
]

\node[stage, fill=evogreen!10]  (raw)      at (0,0)    {原始资料\\Obsidian /\\Notion / MD};
\node[stage, fill=evoyellow!10] (prep)     at (4.2,0)  {预处理\\frontmatter\\atomize / link};
\node[stage, fill=linkblue!10]  (proposal) at (8.4,0)  {ops/proposals/\\提案池};
\node[stage, fill=evored!10]    (curate)   at (12.6,0) {人审 curate\\playbook / facts};

\draw[arrow] (raw) -- (prep);
\draw[arrow] (prep) -- (proposal);
\draw[arrow] (proposal) -- (curate);

% 边标签放在节点上沿之上（y=0.9 > 节点半高 ~0.63），不与节点争夺横向空档。
\node[annot] at (2.1,0.9)  {脚本/LLM};
\node[annot] at (6.3,0.9)  {生成提案};
\node[annot] at (10.5,0.9) {人工};

\end{tikzpicture}
\caption{外部资料导入的两阶段模型}
\label{fig:import}
\end{figure}

\textbf{第一阶段：机器预处理}（可自动化）

\begin{itemize}
  \item \textbf{格式规范化}：补全 YAML frontmatter（id、type、status、triggers、domains、scope），统一文件名；
  \item \textbf{原子化}：把长文档拆成一条一条原子经验（一个知识点/一个检查清单/一个案例一条）；
  \item \textbf{链接转换}：把 wikilink、Notion link 转成 Evo-Kernel 的 id 引用或普通 Markdown link；
  \item \textbf{元数据提取}：用 LLM 或规则自动生成 triggers、domains、scope；
  \item \textbf{暂存提案池}：所有预处理结果放到 \texttt{ops/proposals/import-<批次>/}，不进注入通道。
\end{itemize}

\textbf{第二阶段：人工审查入库}（不可跳过）

\begin{itemize}
  \item 逐条阅读提案，删除噪音、重复、过时的资料；
  \item 把认可的内容改写成 Evo-Kernel 的陈述式经验；
  \item 运行 \texttt{evo curate} 进入 \texttt{facts/}、\texttt{playbook/} 或 \texttt{episodes/}。
\end{itemize}

\begin{redbox}
\textbf{为什么不要一轮直接导入？}
\begin{enumerate}
  \item 绕过 curate 的敏感信息检查，可能把凭据、内网地址、客户信息推入 git；
  \item 绕过指令样审查，可能把祈使式指令（``总是执行 X''）固化成持久化提示注入；
  \item 未经原子化的长文档会降低 recall 精度，大段无关文本稀释 triggers；
  \item 噪音和重复条目直接进入 \texttt{playbook/}，会污染注入集、拉低 helpful/harmful 计数质量。
\end{enumerate}
\end{redbox}

\subsubsection{技术资料进哪个区}

根据资料性质选择目标目录：

\begin{table}[H]
\centering
\caption{技术资料的目标分区}
\begin{tabular}{@{}lp{4.5cm}p{5.5cm}@{}}
\toprule
\textbf{目录} & \textbf{适合内容} & \textbf{示例} \\
\midrule
\texttt{facts/} & 知识点、API 行为、环境配置 & ``Next.js 15 Server Actions 必须 async'' \\
\texttt{playbook/} & 可复用步骤、检查清单、最佳实践 & ``排查 TypeScript 类型错误三步法'' \\
\texttt{episodes/} & 某次具体问题的完整解决过程 & ``2025-07 部署失败排查实录'' \\
\texttt{principles/} & 跨项目/跨技术的通用原则 & ``优先显式类型而非 any'' \\
\bottomrule
\end{tabular}
\end{table}

\subsubsection{示例：从 Markdown 仓库批量导入}

假设你有一个 \texttt{~/Notes/} 目录，里面是一堆 Markdown 技术笔记：

\begin{lstlisting}[language=bash]
# 1. 创建提案池
mkdir -p ~/Dev/evo-kernel/ops/proposals/import-2026-07-24

# 2. 用脚本预处理（这里仅示意，实际根据你的笔记结构写）
python3 scripts/import-notes.py \
  --source ~/Notes \
  --target ~/Dev/evo-kernel/ops/proposals/import-2026-07-24 \
  --format evo-kernel
\end{lstlisting}

预处理后的一个提案文件示例：

\begin{lstlisting}[language=Markdown]
---
id: ts-type-error-three-steps
name: TypeScript 类型错误排查三步法
type: playbook
status: candidate
scope: [debugging]
domains: [typescript]
triggers:
  - TypeScript 类型错误
  - type error
  - 类型排查
evidence:
  helpful: 0
  harmful: 0
verified_by: human
---

遇到 TypeScript 类型错误时，按以下顺序排查：

1. 先看错误位置的上下文，确认是「声明错误」还是「使用错误」；
2. 检查最近的类型断言、泛型参数、`as` 和 `satisfies`；
3. 若仍无法定位，临时用 `// @ts-expect-error` 标注并附带原因，而不是 `any`。
\end{lstlisting}

\begin{lstlisting}[language=bash]
# 3. 人审并 curate（逐条或分批）
cd ~/Dev/evo-kernel
./bin/evo curate --file ops/proposals/import-2026-07-24/ts-type-error-three-steps.md --to playbook
\end{lstlisting}

\subsubsection{研究与借鉴}

Evo-Kernel 的导入模型参考了两种成熟方案：

\begin{itemize}
  \item \textbf{Zettelkasten}：强调 ``fleeting note → literature note → permanent note'' 的多阶段处理。临时捕获的笔记应在 24--48 小时内处理，大部分被丢弃，只有经过自己语言重写的原子思想才进入永久库。
  \item \textbf{RAG ingestion pipeline}：典型的生产级流程包括 extraction → cleaning → chunking → metadata enrichment → embedding → indexing。其中最影响检索质量的不是 chunk 大小，而是 chunk 的语义完整性、上下文线索和元数据治理。
\end{itemize}

这两种方案的共同点是：\textbf{不直接吞入原始资料}，而是经过一个「消化层」——Zettelkasten 靠人工重写，RAG 靠机器分段+元数据+人工校验。Evo-Kernel 的 \texttt{ops/proposals/} + \texttt{evo curate} 就是这个消化层。

\begin{tipbox}
如果资料量很大，可以先用 LLM 批量生成提案草稿，但 \texttt{evo curate} 这一步不要批量自动执行。把 curate 当作 ``最终发布按钮''，保证进入注入通道的每一条都经过人眼确认。
\end{tipbox}

\subsection{工作流五：硬约束硬化}

当某条经验涉及「会造成实际损失的危险操作」，且 skill/提示层已失效时，可以升为 hook。

\begin{lstlisting}[language=bash]
./bin/evo solidify --id no-rm-rf-root --to hook \\
  --matcher "rm\\s+-rf\\s+/" \\
  --match-on command \\
  --mode warn \\
  --message "检测到针对根目录的 rm -rf，请人工确认" \\
  --criteria-confirmed
\end{lstlisting}

\begin{redbox}
\textbf{hook 准入四条件（缺一不可）：}
\begin{enumerate}
  \item \textbf{可程序化判定}：matcher 写得出来、无歧义；
  \item \textbf{违反代价高昂}：不是风格问题，是会造成实际损失的操作；
  \item \textbf{软层已失效}：skill/提示层存在但仍被反复违反（有 harmful 记录佐证）；
  \item \textbf{matcher 覆盖率经复核}：明列覆盖/未覆盖的形态与误报语境。
\end{enumerate}
新 hook 必须先 \texttt{warn} 运行一个 reflect 周期，评估误报率后才能升 \texttt{block}。
\end{redbox}

%======================================================================
\section{周期维护}
%======================================================================

\subsection{每周/每两周做一次}

\begin{lstlisting}[language=bash]
cd ~/Dev/evo-kernel
./bin/evo index rebuild   # 如果 curate/solidify 已经自动 rebuild，可跳过
./bin/evo audit           # 检查治理问题
./bin/evo reflect --save  # 生成周期报告
\end{lstlisting}

\subsection{对账与计数}

helpful/harmful 只能由离线对账单点写（不变量 I4）。主要通道是：

\begin{itemize}
  \item \textbf{自动}：reflect/distill 过程中，Reflector 对每个 session 的注入清单判定四态（adopted / relevant-unused / irrelevant / misleading），写入 \texttt{ops/log/reconcile.jsonl}。
  \item \textbf{人工兜底}：如果你明确知道某条经验被采用或误导了，可以手动标记：
\end{itemize}

\begin{lstlisting}[language=bash]
./bin/evo adopt --ids avoid-silent-bypass        # helpful +1
./bin/evo reject --ids outdated-command-pattern  # harmful +1
\end{lstlisting}

之后需要 \texttt{evo index rebuild} 把计数汇总进 manifest。

\subsection{治理动作}

audit 发现的问题通常对应以下动作：

\begin{table}[H]
\centering
\caption{audit findings 与对应动作}
\begin{tabular}{@{}lp{3.8cm}p{6.2cm}@{}}
\toprule
\textbf{级别} & \textbf{含义} & \textbf{动作} \\
\midrule
HIGH & harmful > helpful & \texttt{evo reject --ids <id>} 或 \texttt{evo demote --id <id> --to archive} \\
MID & superseded 但未归档 & \texttt{evo demote --id <id> --to archive} \\
MID & lessons 中 status 非 candidate & 迁出到正确 zone \\
MID & lessons 条目超过 30 天 & 决定 curate 或 archive \\
LOW & playbook / facts / principles 超过 90 天 & 复验内容是否仍有效 \\
\bottomrule
\end{tabular}
\end{table}

%======================================================================
\section{硬约束执行原理}
%======================================================================

\subsection{guard 与 hook-guard}

\texttt{guard} 是 Pi 的同步调用路径；\texttt{hook-guard} 是 Claude Code PreToolUse hook 的封装。二者共用 \texttt{ops/constraints/*.json} 中的规则。

图 \ref{fig:guard} 展示了 guard 的决策流程。

\begin{figure}[H]
\centering
\begin{tikzpicture}[
  node distance=0.8cm,
  startstop/.style={rectangle, rounded corners, minimum width=2.5cm, minimum height=0.8cm, text centered, draw, fill=evogray!20},
  process/.style={rectangle, rounded corners, minimum width=4cm, minimum height=0.8cm, text centered, align=center, draw, fill=linkblue!10},
  decision/.style={diamond, aspect=2, draw, fill=evoyellow!10, align=center, font=\small},
  arrow/.style={-{Stealth[length=2mm]}, thick}
]

\node[startstop] (start) {工具调用发生};
\node[process, below=of start] (load) {加载 ops/constraints/*.json};
\node[decision, below=of load] (match) {有规则命中？};
\node[decision, below left=0.6cm and 1cm of match] (mode) {mode = block？};
\node[process, below=of mode] (deny) {deny\\阻断调用};
\node[process, right=2.5cm of mode] (warn) {warn\\浮现提示但放行};
\node[process, right=2.5cm of match] (allow) {allow\\静默放行};

\draw[arrow] (start) -- (load);
\draw[arrow] (load) -- (match);
\draw[arrow] (match) -- node[left, font=\scriptsize] {命中} (mode);
\draw[arrow] (match) -- node[above, font=\scriptsize] {未命中} (allow);
\draw[arrow] (mode) -- node[left, font=\scriptsize] {是} (deny);
\draw[arrow] (mode) -- node[above, font=\scriptsize] {否} (warn);

\end{tikzpicture}
\caption{guard 决策流程}
\label{fig:guard}
\end{figure}

\begin{warnbox}
\textbf{fail-open 原则：} guard 自身异常（约束文件损坏、正则非法等）时，一律按 allow 处理。只有「规则正确匹配且 mode=block」才会 deny。这是为了保证经验治理层不会反过来阻塞正常工作。
\end{warnbox}

%======================================================================
\section{备份与恢复}
%======================================================================

\subsection{备份机制}

Evo-Kernel 有两层备份：

\begin{enumerate}
  \item \textbf{主备份}：每次 \texttt{curate} / \texttt{solidify} / \texttt{demote} 自动 \texttt{git push} 到私有 remote（RPO = 一次提交）。
  \item \textbf{兜底备份}：定期生成 git bundle 到本地异地介质：
\end{enumerate}

\begin{lstlisting}[language=bash]
cd ~/Dev/evo-kernel
git bundle create ~/evo-kernel-$(date +%Y%m%d).bundle --all
\end{lstlisting}

\begin{redbox}
\textbf{原始日志不备份：} \texttt{ops/log/*.jsonl} 含 prompt 明文，已 \texttt{gitignore}，永不入 remote/bundle。盘损时原始判据数据会丢失，只有 reflect 聚合后的历史保留。这是 §4.4 的诚实权衡。
\end{redbox}

\subsection{恢复演练}

每季度做一次：

\begin{lstlisting}[language=bash]
TMPDIR=$(mktemp -d)
cd "$TMPDIR"
git clone git@github.com:zodyne/evo-kernel.git
cd evo-kernel
npm install
./bin/evo doctor   # 此时 harness 挂线项会 WARN/FAIL，属预期
./test/smoke.sh    # 必须 PASS=85 FAIL=0
rm -rf "$TMPDIR"
\end{lstlisting}

%======================================================================
\section{红线与注意事项}
%======================================================================

\begin{redbox}
\begin{enumerate}
  \item \textbf{remote 必须私有。} 经验条目含工作语境，公开 remote 直接违反 §4.3 红线。
  \item \textbf{inbox / lessons 永不注入。} 只有 curate 后的分区才有资格进入 recall 通道（不变量 I2）。
  \item \textbf{helpful/harmful 只由 reconcile 单点写。} 禁止实时反馈回路（不变量 I4）。
  \item \textbf{curate 必须人审。} 这是唯一入库口，负脱敏与指令样审查责任（不变量 I3）。
  \item \textbf{原始日志不入 git。} \texttt{ops/log/*.jsonl} 全部 \texttt{gitignore}（不变量 I4 / §4.4）。
  \item \textbf{hook 升 block 前必须过 warn 观察期。} 且需人工复核 matcher 覆盖率（§8）。
  \item \textbf{不要多机并写。} remote 仅备份，第二台机器只读消费；多机并写会 ID 碰撞、计数合并未定义（§4.1）。
\end{enumerate}
\end{redbox}

\subsection{常见误区}

\begin{itemize}
  \item \textbf{「我把文件拖进 playbook/ 就行了吧？」} 不行。curate 是唯一能写入 playbook/facts/episodes/principles 的入口。
  \item \textbf{「recall 空命中很多，是不是该上 M1 FTS5？」} M1 判据是「注入精度连续 2 周期 <50\% 且 recall $\geq$ 200 次」。截至 2026-07-29 recall 累计 193 次、召回精度 54\%，两条都未命中；且对账覆盖率仅 20\%，精度本身还不可解读。做了也是过早优化。
  \item \textbf{「transcript 哨兵比高，是不是 bug？」} 不是 bug。Claude Code 会按策略清理 transcript 文件，登记与蒸馏之间存在腐烂窗口。由 reflect 周期跟踪，不是靠实施解决。
  \item \textbf{「我想在另一台机器上也 curate。」} 暂时不要。多机并写未设计，只读消费可以。
\end{itemize}

%======================================================================
\section{实施进度评估（截至 2026-07-29）}
%======================================================================

\begin{tipbox}
本节所有数字取自 2026-07-29 当日在本机实跑的 \texttt{evo doctor --full} 与 \texttt{evo reflect} 输出，非估算。复现方式：\texttt{cd \textasciitilde/Dev/evo-kernel \&\& ./bin/evo doctor --full \&\& ./bin/evo reflect}。

本节数字的时效性本身就是一条经验（条目 \texttt{doc-selfreported-counts-drift}）：初稿按 07-28 数据写成，隔一天复核时库规模已从 50 涨到 67、smoke 从 85 涨到 93、待审提案从 11 涨到 22。**读到本节请先按上面的命令复跑一遍**，不要直接引用纸面数字。
\end{tipbox}

\subsection{总体结论}

内核已完整落地并进入日常生产使用：命令面 26 个、smoke 93 项断言全绿、双 harness 挂线完好、飞轮（capture $\rightarrow$ recall $\rightarrow$ 蒸馏 $\rightarrow$ curate）已闭环并自动化运行。设计文档 \texttt{design/final-architecture.md} §9.3 列出的 4 项「现在就做」（无判据）清单已全部核销。

Phase 1--3（FTS5 / 向量 / LightRAG）按设计全部处于判据门控状态，当前无一命中——这是符合演进纪律的预期状态，不是欠账。

真正的短板不在实现，在**治理链路的下游**：固化梯度（solidify）与退役出口（demote）代码就绪但零生产实例，硬约束 2 条规则全部停在 warn 且误报率 55--71\%，离线对账覆盖率仅 20\%（低于 30\% 的可解读门槛）。

\subsection{Phase 0 清单核销}

\begin{table}[H]
\centering
\caption{design/final-architecture.md §9.3「现在就做」四项}
\begin{tabular}{@{}lp{5.5cm}p{4.2cm}@{}}
\toprule
\textbf{项} & \textbf{状态} & \textbf{证据} \\
\midrule
M0 解析器修复 & \textcolor{evogreen}{已完成} & 引入 \texttt{js-yaml}，且显式用 \texttt{CORE\_SCHEMA} 堵住日期被静默转 \texttt{Date} 的坑（commit \texttt{ed0c80e}） \\
K0a primer & \textcolor{evogreen}{已完成} & \texttt{evo primer --install}；doctor 项 17 PASS：两处已装，\texttt{review\_after} 2026-09-25 \\
去锚定 L1 & \textcolor{evogreen}{已完成} & primer 常驻块落 \texttt{CLAUDE.md}/\texttt{AGENTS.md} 为主、kernel 注入为辅（commit \texttt{312e62c}） \\
日常飞轮 & \textcolor{evogreen}{运行中} & recall 累计 193 次；库从 18 条增至 67 条；蒸馏已自动化 \\
\bottomrule
\end{tabular}
\end{table}

\subsection{能力落地矩阵}

\begin{table}[H]
\centering
\caption{七条主链路的实现与验证状态}
\begin{tabular}{@{}lp{2.4cm}p{6.4cm}@{}}
\toprule
\textbf{能力} & \textbf{状态} & \textbf{实测证据} \\
\midrule
捕获 capture & \textcolor{evogreen}{生产可用} & inbox 零判断入口在用；积压已按「junk drawer 前兆」清理过一轮（\texttt{420ec73}） \\
检索注入 recall & \textcolor{evogreen}{生产可用} & 193 次调用，76\% 有注入命中（147/193）；后端仍为 \texttt{scan}（Phase 1 未触发） \\
离线蒸馏 distill & \textcolor{evogreen}{已自动化} & SessionEnd 触发器 + launchd 每日 03:30 兜底（\texttt{--max 3}、带锁）；\texttt{queue} / \texttt{mark-distilled} / \texttt{catalog} 三命令支撑队列与查重 \\
人审入库 curate & \textcolor{evogreen}{生产可用} & 51 条 \texttt{validated} 已入库；入 validated 区改写 status 的状态机漏洞已修（\texttt{cc45bf6}）；未跟踪提案致 \texttt{git add} 整批 fatal、commit 从未成功的缺陷已修（\texttt{cc3ff21}） \\
离线对账 reconcile & \textcolor{evoyellow!70!black}{部分验证} & 39 例已判四态（adopted 6 / relevant-unused 15 / irrelevant 18 / misleading 0），覆盖率仅 20\% \\
技能固化 solidify & \textcolor{evored}{路径未验证} & 代码与 smoke 就绪，但\textbf{生产实例为 0}；现有 4 个 SKILL 均为手写运维技能，非条目固化产物；reflect 已挂出 6 条固化候选待处理 \\
硬约束 guard & \textcolor{evored}{路径未验证} & 2 条规则，\textbf{均停在 warn}：\texttt{dangerous-rm-rf} 误报率 71\%（59 裸命中/42 引号内提及）、\texttt{evo-curate-needs-human} 误报率 55\% 且 \texttt{criteria\_confirmed=false}，双双未达 §8 准入第 4 条 \\
\bottomrule
\end{tabular}
\end{table}

\subsection{库存与飞轮实测数据}

\begin{table}[H]
\centering
\caption{2026-07-29 实测规模}
\begin{tabular}{@{}lp{4.6cm}p{4.6cm}@{}}
\toprule
\textbf{指标} & \textbf{当前值} & \textbf{备注} \\
\midrule
条目总数 & 67 & 其中可注入 51 条（\texttt{validated}） \\
分区分布 & playbook 41 / lessons 9 / inbox 7 / facts 5 / episodes 4 / principles 1 & principles 单薄是真实缺口 \\
待审提案 & 22 & 一天内从 11 翻倍，后台蒸馏产出快于人审 \\
\texttt{ops/archive/} & 0 & 退役出口从未被走过 \\
skills & 4 & 均为手写运维技能 \\
constraints & 2（均 warn） & 无 block 规则 \\
CLI 命令面 & 26 & 设计资产表记录的 617 行已增至 1463 行 \\
smoke 断言 & PASS=93 FAIL=0 & 手册旧版记录为 61 \\
doctor & 16 PASS / 1 WARN / 0 FAIL & WARN = 工作区有未提交变更 \\
recall 累计 & 193 次 & 空命中任务 46 个 \\
transcript 哨兵比 & 36/81（44\%） & 已由 FAIL/WARN 转为 PASS \\
\bottomrule
\end{tabular}
\end{table}

\subsection{判据门控状态}

\begin{table}[H]
\centering
\caption{Phase 1--3 判据（全部未命中，符合预期）}
\begin{tabular}{@{}lp{3.1cm}p{3.6cm}p{2.9cm}@{}}
\toprule
\textbf{判据} & \textbf{阈值} & \textbf{当前值} & \textbf{结论} \\
\midrule
M1 recall 累计 & $\geq$ 200 & 193 & 未命中（已逼近） \\
M1 召回精度（检索层） & 连续 2 周期 $<$ 50\% & 54\%（21/39） & 未命中 \\
采纳率（应用层） & 观察项，阈值待数据 & 29\%（6/21） & 观察 \\
M1 对账覆盖率 & $\geq$ 30\% 方可解读 & 20\%（39/193） & \textcolor{evored}{不足，精度暂不可解读} \\
M2 条目数（向量） & $\geq$ 500 & 67 & 未命中 \\
P4 多跳关系需求 & 空命中呈关系特征 & 46 个空命中，多为闲聊/短指令 & 未命中 \\
\bottomrule
\end{tabular}
\end{table}

\begin{warnbox}
对账覆盖率 20\% 是当前最关键的观测缺口：精度建立在 193 次注入里仅 39 次被判定的基础上，样本偏斜风险高。在覆盖率提到 30\% 以前，不应据此讨论换检索后端。

另需注意精度口径已拆分为两层：\textbf{召回精度 54\%}（检索层——找出来的条目是否相关）与 \textbf{采纳率 29\%}（应用层——相关条目是否真被用上）。两者差距说明当前瓶颈更可能在「注入了但没被采纳」，而不在检索本身——这进一步支持「先修蒸馏与采纳纪律，而非先换检索后端」。
\end{warnbox}

\subsection{尚未验证的路径（诚实清单）}

\begin{enumerate}
  \item {\sloppy \textbf{固化梯度是空转的。} \texttt{solidify --to skill} 从未在生产跑过一次，\texttt{ops/archive/} 为空。这意味着「经验 $\rightarrow$ 技能」和「经验退役」这两条设计承诺的出口，只有 smoke 覆盖、没有真实使用反馈。reflect 当前挂着 6 条固化候选（最高的是 \texttt{verify-external-references}，helpful=5）与 1 条退役候选，是消除该空白的现成入口。\par}
  \item \textbf{硬约束层实质未启用。} 2 条规则全部停在 \texttt{warn}：matcher 是子串匹配，分不清「执行」与「引号内提及」，误报率分别为 71\% 与 55\%，reflect 对两条都明确不建议升 \texttt{block}。硬化梯度目前仍是纸面能力。
  \item \textbf{principles 仅 1 条。} 跨域普适原则本应是最高价值层，实际产出集中在 playbook（41 条）。蒸馏产物的层级分布偏低。
  \item \textbf{提案积压 22 条。} 人审是设计上的瓶颈也是必要环节，但积压量一天内从 11 翻倍到 22——后台蒸馏的产出速度已明显快于人审速度，「消化层反成堆积层」不再是风险而是正在发生。
  \item \textbf{多机并写仍未设计。} remote 只作备份，第二台机器只读消费的约束未变。
\end{enumerate}

\subsection{建议的下一步（按优先级）}

\begin{enumerate}
  \item 走通一次完整的 \texttt{solidify --to skill}，用 reflect 挂出的高 helpful 候选打通固化出口；同步走一次 \texttt{demote} 清掉空转条目。
  \item 提升对账覆盖率至 30\% 以上——修蒸馏与采纳纪律优先于修检索（采纳率 29\% 明显低于召回精度 54\%）。
  \item 收窄两条约束的 matcher（剥离引号语境，区分执行与提及），再谈 warn $\rightarrow$ block。
  \item 消化 22 条待审提案，并给后台蒸馏加产出节流——否则人审永远追不上。
  \item Phase 1--3 继续保持门控，不提前动工。
\end{enumerate}

%======================================================================
\section{命令速查表}
%======================================================================

\begin{table}[H]
\centering
\caption{Evo-Kernel 26 命令速查}
\begin{tabular}{@{}lll@{}}
\toprule
\textbf{命令} & \textbf{用途} & \textbf{典型场景} \\
\midrule
\texttt{capture "..."} & 零判断捕获 & 随时记录想法 \\
\texttt{recall --task "..."} & 主动检索 & 手动查经验 \\
\texttt{hook-recall} & harness 注入入口 & 由 hooks 自动调用 \\
\texttt{session-end} & 手动登记会话 & hook 未触发时 \\
\texttt{hook-session-end} & 自动登记入口 & 由 hooks 自动调用 \\
\texttt{adopt --ids} & 人工 +helpful & 明确知道某条被采用 \\
\texttt{reject --ids} & 人工 +harmful & 明确知道某条误导 \\
\texttt{reconcile --ids --state} & Reflector 四态对账 & 蒸馏时自动写 \\
\texttt{queue} & 待蒸馏队列 TSV & 后台驱动器输入 \\
\texttt{mark-distilled --ids} & 回写蒸馏完成 & 蒸馏收尾 \\
\texttt{curate --file --to} & 人审入库 & 核心流程 \\
\texttt{slice --session} & transcript 切片 & 蒸馏前准备 \\
\texttt{index rebuild} & 重建 manifest & curate 后自动做 \\
\texttt{primer [--install]} & 常驻背景块 & 装到 CLAUDE.md / AGENTS.md \\
\texttt{candidates} & 列出候选条目 & agentic 粗筛 \\
\texttt{catalog} & id×triggers 查重清单 & Reflector 落笔前查重 \\
\texttt{get --ids} & 取全文 & agentic 精筛 \\
\texttt{reflect --save} & 周期复盘 & 每周/每两周 \\
\texttt{audit} & 治理检查 & 配合 reflect \\
\texttt{inbox} & 待处理清单 & 查看 capture 和 refs \\
\texttt{solidify --id --to skill} & 技能固化 & 经验变 SKILL.md \\
\texttt{solidify --id --to hook} & 硬约束注册 & 危险操作防护 \\
\texttt{demote --id --to archive} & 退役 & 经验过期/有害 \\
\texttt{guard / hook-guard} & 约束执行 & 由 harness 调用 \\
\texttt{link} & 同步 skills 软链 & solidify 后自动做 \\
\texttt{doctor [--full]} & 部署自检 & 每日/每次重大变更后 \\
\bottomrule
\end{tabular}
\end{table}

\end{document}
